\pdfvariable trailerid {[<C3442026090300000000000000000000><C3442026090300000000000000000000>]}
\documentclass[10pt]{article}
\usepackage[margin=0.72in]{geometry}
\usepackage{amsmath,amsthm,unicode-math,microtype,xurl}
\providecommand{\CRevisionRound}{2}
\newtheorem{theorem}{Theorem}
\newtheorem{remark}[theorem]{Remark}
\newcommand{\CC}{\mathbb C}
\newcommand{\RR}{\mathbb R}
\newcommand{\sn}{\operatorname{sn}}
\newcommand{\cn}{\operatorname{cn}}
\newcommand{\dn}{\operatorname{dn}}
\newcommand{\PiK}{\Pi}
\setlength{\emergencystretch}{3em}
\title{Complete Elliptic Reduction and Two-Phase Return\\for a Hamiltonian Resonant Triad}
\author{Route-A source-local certificate HCS-C344}
\date{3 September 2026}

\begin{document}
\maketitle
\begin{abstract}
For the canonical complex resonant triad we prove global Liouville
integrability and reduce every intensity orbit to a real cubic.  Every
regular orbit is written in Jacobi form with its exact intensity period.
\ifnum\CRevisionRound>0
Two complete elliptic integrals of the third kind reconstruct the surviving
torus phases and give the necessary and sufficient full-state return test.
The zero-Hamiltonian transfer, its factor-two period, the equal-action
separatrix, and maximal-Hamiltonian relative equilibria are closed separately.
\fi
\ifnum\CRevisionRound>1
Finite exact receipts are separated from the continuum proof, and source,
collision, and Route-A boundaries are recorded explicitly.
\fi
\end{abstract}

\section{Canonical convention and first integrals}
On $\CC^3$ use
\[
 \{f,g\}=-i\sum_{j=1}^3
 \left(f_{z_j}g_{\bar z_j}-f_{\bar z_j}g_{z_j}\right)
\]
and the real Hamiltonian
\begin{equation}\label{eq:H}
 H=z_1z_2\bar z_3+\bar z_1\bar z_2z_3.
\end{equation}
Hamilton's equations are
\begin{equation}\label{eq:triad}
 i\dot z_1=\bar z_2z_3,\qquad
 i\dot z_2=\bar z_1z_3,\qquad
 i\dot z_3=z_1z_2.
\end{equation}
Writing $I_j=|z_j|^2$ and $w=z_1z_2\bar z_3$ gives
\[
 \dot I_1=\dot I_2=-2\operatorname{Im}w,\qquad
 \dot I_3=2\operatorname{Im}w.
\]
Consequently
\begin{equation}\label{eq:MR}
 N_1=I_1+I_3,\qquad N_2=I_2+I_3,\qquad H
\end{equation}
are conserved.  Direct evaluation of the displayed bracket gives pairwise
involution.  On the open set $I_1I_2I_3\sin\Psi\ne0$, where
$\Psi=\arg z_1+\arg z_2-\arg z_3$, the first two differentials are
independent action differentials and
\[
 \partial_\Psi H=-2\sqrt{I_1I_2I_3}\sin\Psi\ne0.
\]
Thus the three integrals are generically independent.  Their common levels
are bounded because $0\le I_3\le\min(N_1,N_2)$ and
$I_j=N_j-I_3$ for $j=1,2$.  Polynomial local existence and boundedness give
a global flow on all of $\RR$.

\section{The real cubic and every regular intensity}
Put $x=I_3$.  Since $\dot x=2\operatorname{Im}w$ and
$|w|^2=x(N_1-x)(N_2-x)$,
\begin{equation}\label{eq:cubic}
 \dot x^2=4x(N_1-x)(N_2-x)-H^2=:P(x).
\end{equation}
Let $N_-:=\min(N_1,N_2)$ and $N_+:=\max(N_1,N_2)$.  The function
$f(x)=x(N_1-x)(N_2-x)$ has one maximum on $[0,N_-]$, at
\begin{equation}\label{eq:xstar}
 x_*=\frac{N_1+N_2-\sqrt{N_1^2-N_1N_2+N_2^2}}{3},\qquad
 H_{\max}=2\sqrt{f(x_*)}.
\end{equation}
For $0<|H|<H_{\max}$, signs at $0$, the maximum, $N_-$, $N_+$,
and infinity prove that $P$ has exactly three roots
\begin{equation}\label{eq:rootorder}
 0<r_1<r_2<N_-\le N_+<r_3.
\end{equation}
They obey
\begin{equation}\label{eq:vieta}
 \sum r_j=N_1+N_2,\quad
 \sum_{j<k}r_jr_k=N_1N_2,\quad
 r_1r_2r_3=\frac{H^2}{4}.
\end{equation}

\begin{theorem}[regular elliptic orbit]\label{thm:regular}
Let $0<|H|<H_{\max}$ and define
\[
 m=\frac{r_2-r_1}{r_3-r_1},\qquad
 u=\sqrt{r_3-r_1}(t-t_0).
\]
Every intensity orbit on this level is
\begin{equation}\label{eq:snsolution}
 x(t)=r_1+(r_2-r_1)\sn^2(u\mid m),
 \qquad T_x=\frac{2K(m)}{\sqrt{r_3-r_1}}.
\end{equation}
\end{theorem}

\begin{proof}
On $[r_1,r_2]$ the cubic is
$4(x-r_1)(r_2-x)(r_3-x)$.  Substitute~\eqref{eq:snsolution} and use
$(\sn')^2=(1-\sn^2)(1-m\sn^2)$.  The resulting identity is
exact.  Time translation selects either initial sign of $\dot x$, and
$\sn^2$ has least real period $2K(m)$.
\end{proof}

The analytic root argument and substitution are the \emph{cubic elliptic
owner} of revision round zero.

\ifnum\CRevisionRound>0
\section{Two phases, not one, control full return}
When $H\ne0$, no $I_j$ can vanish.  Writing
$z_j=\sqrt{I_j}e^{i\phi_j}$, division of~\eqref{eq:triad} by $z_j$ gives
\begin{equation}\label{eq:phasedot}
 \dot\phi_1=-\frac{H}{2(N_1-x)},\qquad
 \dot\phi_2=-\frac{H}{2(N_2-x)},\qquad
 \dot\phi_3=-\frac{H}{2x}.
\end{equation}
With $\PiK(n\mid m)$ denoting the complete third-kind integral, integration
over one intensity period yields
\begin{equation}\label{eq:deltas}
 \Delta_j=-\frac{H}{\sqrt{r_3-r_1}(N_j-r_1)}
 \PiK\left(\frac{r_2-r_1}{N_j-r_1}\mathrel|m\right),
 \qquad j=1,2.
\end{equation}
Indeed, the integral over $0\le u\le2K$ is twice its complete-quarter
integral.  After $T_x$, both $x$ and $\dot x$ return; together with $H$
they determine $e^{i\Psi}$.  Hence
$\Delta_3=\Delta_1+\Delta_2$ modulo $2\pi$.

Because a return of the nonconstant full state must occur at an integer
multiple of its least intensity period, the full complex orbit is periodic
if and only if
\begin{equation}\label{eq:closure}
 \frac{\Delta_1}{2\pi}\in\mathbb Q,\qquad
 \frac{\Delta_2}{2\pi}\in\mathbb Q.
\end{equation}
One scalar period therefore does not settle full-state recurrence.

\section{Zero Hamiltonian and double-root boundaries}
Suppose first that $0<N_1<N_2$.  With $A=\sqrt{N_1}$,
$B=\sqrt{N_2}$, $m=A^2/B^2$, and $u=B(t-t_0)$, every non-equilibrium
orbit, up to the two phase symmetries, is
\begin{equation}\label{eq:Hzero}
 \begin{aligned}
 z_1&=A\cn(u\mid m)e^{i\alpha},&
 z_2&=B\dn(u\mid m)e^{i\beta},\\
 z_3&=-iA\sn(u\mid m)e^{i(\alpha+\beta)}.
 \end{aligned}
\end{equation}
The other ordering exchanges modes $1$ and $2$.  The Jacobi derivative
identities verify~\eqref{eq:triad} directly, including both amplitude-zero
crossings.  Intensities have period $2K(m)/B$.  Translation by $2K$
changes the signs of $\cn$ and $\sn$ but not $\dn$, so the full state has
least period $4K(m)/B$.

At $N_1=N_2=N>0$, the limit $m=1$ gives
\begin{equation}\label{eq:separatrix}
 \sqrt N\left(\operatorname{sech}u\,e^{i\alpha},
 \operatorname{sech}u\,e^{i\beta},
 -i\tanh u\,e^{i(\alpha+\beta)}\right),\qquad
 u=\sqrt N(t-t_0).
\end{equation}
It is heteroclinic between opposite points on the $z_3$-axis and has
infinite period.

At $|H|=H_{\max}$, the two accessible roots coalesce at $x_*$.  The
intensities are constant and the phases have frequencies
\begin{equation}\label{eq:omegas}
 \omega_1=-\frac{H}{2(N_1-x_*)},\quad
 \omega_2=-\frac{H}{2(N_2-x_*)},\quad
 \omega_3=-\frac{H}{2x_*}=\omega_1+\omega_2.
\end{equation}
The final equality is equivalent to $f'(x_*)=0$.  This relative equilibrium
is periodic exactly when $\omega_1/\omega_2$ is rational.  The symmetric
face is always periodic; the asymmetric example $(N_1,N_2)=(5,8)$ has
$x_*=2$, $H_{\max}=12$, and absolute frequencies $(2,1,3)$.

The origin and each complex coordinate axis are equilibrium families.
Introducing a real coupling $g$ multiplies the vector field by $g$: the
$g=0$ face is the identity flow, and $g\ne0$ is the present theorem after
$\tau=gt$.  Complex conjugation paired with time reversal leaves all
intensities unchanged.  These clauses form the \emph{two-phase return owner}
added in revision round one.
\fi

\ifnum\CRevisionRound>1
\section{Executable receipts, collisions, and Route A}
The evidence ledger contains $72$ regular signed-Hamiltonian rows, $12$
zero-Hamiltonian rows, and $24$ relative-equilibrium rows.  Every regular
root has an exact rational sign bracket.  A producer-independent checker
recomputes the roots, $K$, both complete third-kind integrals, the factor-two
boundary period, and every relative frequency.  A separate symbolic lane
checks the Poisson signs, involution, cubic, Jacobi substitution, and explicit
zero-Hamiltonian solution.  Sampling proves none of the continuum claims;
the preceding arguments do.

The nearest workspace systems are different.  C211 treats a Hamiltonian
Lotka--Volterra period annulus, C230 an open Toda lattice, C235 cyclic
population dynamics, and C256 a KdV traveling-wave profile.  None owns a
complex resonant triad with two surviving phase returns.

The conservative Route-A tuple is
\[
 (\mathrm{A0\_FAIL},\mathrm{A1\_WEAK},\mathrm{A2\_FAIL},
  \mathrm{A3\_FAIL},\mathrm{A4\_FORMAL\_HINT}).
\]
Elliptic source dynamics gives only A1 weakly.  A formal three-boson cubic
analogy does not supply a self-adjoint domain theorem or a complete quantum
spectrum.  There is no rational-prime carrier, prime-power clock, arithmetic
orbit dictionary, target Euler product, root number, automorphy, target
divisor or functional equation, target-zero match, or Hilbert--Polya
operator.  Route B is false.  This is the \emph{finite evidence and route
firewall} of the final revision.
\fi

\section*{Source boundary}
This paper is a source-local reconstruction, not a priority claim.
Manley and Rowe own their energy relations; Armstrong et al.\ own the
coupled optical amplitude setting and explicit three-wave solutions; Kaup,
Reiman, and Bers give an authoritative primary treatment of broader
resonant three-wave dynamics.  The exact software receipts are local
implementation checks and are not attributed to those sources.

\begin{thebibliography}{9}
\bibitem{ManleyRowe}
J. M. Manley and H. E. Rowe, ``Some General Properties of Nonlinear
Elements---Part I. General Energy Relations,'' \emph{Proceedings of the IRE}
44 (1956), 904--913, DOI \url{10.1109/JRPROC.1956.275145}.
\bibitem{Armstrong}
J. A. Armstrong, N. Bloembergen, J. Ducuing, and P. S. Pershan,
``Interactions between Light Waves in a Nonlinear Dielectric,''
\emph{Physical Review} 127 (1962), 1918--1939,
DOI \url{10.1103/PhysRev.127.1918}.
\bibitem{Kaup}
D. J. Kaup, A. Reiman, and A. Bers, ``Space-time evolution of nonlinear
three-wave interactions. I. Interaction in a homogeneous medium,''
\emph{Reviews of Modern Physics} 51 (1979), 275--309,
DOI \url{10.1103/RevModPhys.51.275}.
\end{thebibliography}

\end{document}
